
This work investigated whether benchmark design features create systematic constraints enabling predictive coverage analysis. Using a pilot study of 20 diverse benchmarks, design features were extracted objectively (Cohen's kappa $\geq 0.917)$ and clustered into coverage families that predict future citation co-occurrence with 78\% accuracy.

Three main contributions were demonstrated. First, temporal persistence: pre-2023 features predict 2023-2024 patterns with 78.07\% citation overlap, outperforming random (19.61\%) by 585\%. Second, standardized extraction achieving substantial agreement: kappa 0.917-1.0 enables scaling from pilot to large-scale analysis. Third, modality-driven clustering: 0.748 average intra-family similarity reveals modality creates stronger coverage patterns than task formulation.

Future work should address untested alternatives: real-world citation validation on 1000+ ArXiv papers to measure precision degradation from synthetic (100\%) to realistic (expected 75-85\%), and task-based clustering ablation to isolate modality vs task contributions. Scope extensions include scaling to 100+ benchmarks covering rare modalities (video, 3D, tabular) and fine-grained subclusters (k=8-12) to capture task-based patterns within modality groups. Unverified assumptions include cross-temporal robustness (pre-2020 $\rightarrow$ 2024) to assess 4-year persistence, and citation threshold sensitivity analysis to determine minimum coverage requirements.

Beyond scaling, this work enables automated coverage gap discovery---identifying hypothesis categories with zero existing benchmarks before implementation begins---and potential integration into research planning tools to estimate validation feasibility during hypothesis formulation.
